\section{Sustainability}\label{sustainability}

\stanceinfo{\textbf{CFES Congress 2018} [2018-01-07]}{2024-01-02}

\officialstance{The Canadian Federation of Engineering Students believes that sustainability is an essential consideration in engineering practice, and believes that it is necessary to educate and engage its members on issues of sustainability, while also evaluating and improving the sustainability in engineering curriculum.}

\stanceheading{The Students' Position}
\begin{itemize}
 \item Engineers have a responsibility to hold paramount the safety, health and welfare of the public and have regard for the environment, which entails a responsibility to perform work that is sustainable for future generations.
 \item The future generations of engineers should be thoroughly educated and invested in sustainable practices of engineering prior to graduation, through proper applications within the engineering curriculum.
 \item The CFES should evaluate the sustainability of its operations to reform and improve its practices as necessary.
 \item Advocacy work surrounding sustainability should be assessed within the lens of engineering students, job opportunities and the future of engineering.
\end{itemize}

\stanceheading{Recommendations to Partners, Stakeholders, and Other Entities}
\begin{itemize}
 \item The CFES calls on the Canadian Engineering Accreditation Board (CEAB) to include sustainability-focused Accreditation Units (AUs) for all engineering degrees, and to ensure that every engineering student is graduating with those competencies.
 \item The CFES calls on members to perform institution-specific level advocacy focusing on including sustainability considerations in their everyday operations and in travel to respective CFES conferences (i.e. encouraging institutions to divest and provide sustainability resources).
 \item The CFES calls on Engineering Deans Canada (EDC) to support sustainability initiatives, teams and practices within their individual institutions and partner with organizations to create a more sustainable engineering campus.
\end{itemize}

\stanceheading{What the CFES is doing}
\begin{itemize}
 \item The CFES hosts a national conference on sustainability in engineering annually at one of its member schools with the goal of directly educating and engaging engineering students on sustainable practices related to the profession that they can bring back and implement ideas into their universities.
 \item The CFES will continue to update the best practices for sustainability for use by both itself and its member societies Sustainability Bank of Best Practices, as well as other useful learning resources related to sustainability.
\end{itemize}

\stanceheading{What the CFES plans to do}
\begin{itemize}
 \item The CFES will perform an annual environmental footprint analysis of its operations, and reform its operations to align with sustainable practices wherever possible.
 \item The CFES will assist to develop documentation for a conference guideline on sustainability to consider sustainable practices in their planning and execution processes. As a national student organization, the CFES has a major role to play in paving the way for sustainable actions as an organization and then further motivating the member schools.
 \item The CFES can develop guiding material for itself and member organizations to decrease their environmental impact, through a living document outlining best practices. This document shall be reviewed and updated annually to reflect any changes in best practices.
\end{itemize}

\newpage
